\section{Introduction}

Mundelein Center is a tall, narrow classroom building. Its classrooms are spread from floor~1 to floor~14, and
all of them are reached through two elevator banks of four and two cars. Only the two-car bank travels above
floor~9. Every 75 minutes on Tuesdays and Thursdays and every 50 minutes on Mondays, Wednesdays, and
Fridays, students leaving one class and arriving for the next converge on the elevators. Cars leave the upper
floors carrying four to six people and still stop at every floor below, even though they are physically full,
so students on the lower floors wait for cars that cannot take them.

We focus on classes that end before 4~PM and on two behavioral measures that cost nothing to implement: asking
students to take their backpacks off and hold them in front (our estimate: two more riders per car), and asking
students on floors 2 and~3 to walk. A third question is scheduling: because the registrar's grid fixes the
class times, can the upper floors run on a shifted clock so that their peak does not coincide with the rest of
the building?

The contribution is a transparent, reproducible model, not a measurement; the classical treatments of elevator traffic~\cite{barney,strakosch} motivate its structure. All code, parameters, and result
tables accompany this paper (Section~\ref{sec:repro}).

A reader should keep two kinds of uncertainty apart. The first is \emph{parametric}: elevator speed, door time,
boarding time per rider, the number of extra riders gained by holding backpacks, and the compliance rates are
estimates. We can vary them, and we do. The second is \emph{structural}: the simulator encodes a particular
account of how students behave in a queue and how the elevators operate, and that account is not a parameter that
can be swept. For example, the model gives students no way to give up on the line, represents partial compliance
as a coin flip per car, and fills every seat of every class. A structural assumption can change a result as much
as any parameter we vary, but it will not show up in a sensitivity table, because it is the same in every run.
Section~\ref{sec:struct} lists these assumptions, and the results, discussion, and recommendations are written
with them in view.

We also state what kind of contribution this is. The paper presents a simulation framework and preliminary
estimates of the effect of two measures. Its conclusions are model results, that is, what follows under stated
assumptions. Some are robust to the parameters we varied, but none is yet an empirical prediction for the building,
and the paper is not a basis for operational change. Section~\ref{sec:future} ranks the work needed to move from
model results toward evidence.
